% STATUS: PHASE 1 REVIEW COMPLETED WITH EXPLICIT ACCESS AND NOVELTY GAPS
% This is an includable fragment; see literature_audit_compendium.tex.
% No proposed project theorem is proved or promoted by this audit.
\section{Literature audit: scope and evidence}
\label{sec:phase1-audit}

\subsection{Contract, corpus and limits}
This audit addresses the eight novelty questions for the canonical
demand-modulated model $(\DM)_{n,q}$ and its matched bilinear comparator
$(\MB)_{n,q}$. Their equations and assumptions remain exactly those in
\texttt{00\_canonical\_model.tex}. The main scope is arbitrary $n\geq5$ and
$1\leq q\leq n-1$, with the exact selection program concentrated on $q=2$.
Five-element literature is integrated wherever it affects the model,
coupling law, $n=5,q=2$ specialization or comparator; it is not a separate
archival research track.

The six author-supplied PDFs form the core corpus, S01--S06. They were
accessible at the local paths supplied on 20 September 2026, although the
project's synced \texttt{sources/} directory still contained only the master
blueprint. S07--S14 are targeted primary-source supplements or expressly
unresolved leads. No master-blueprint statement is evidence of novelty.
No project theorem has been derived, no experiment run, and no source's
numerical findings independently reproduced.

\paragraph{Locator convention.}
For S01--S04, page numbers below are the \emph{printed pages in the supplied
author PDFs}, which also equal their one-based PDF positions. They are not
inferred journal pages. S05 has printed pages 777--781 (PDF positions 1--5);
S06 has publisher pages 1--15. The local filenames, page counts and SHA-256
digests are in \texttt{audits/phase1\_source\_manifest.md}. Supplemental
versions and pagination are identified individually. A missing full-text
locator is recorded as an access gap, never manufactured from an abstract.

\paragraph{Method.}
The audit compared each displayed source model with the canonical coupling
locations, signs, state dependence, dimension and parameter restrictions.
It inspected the relevant statements and calculations for spectra, mode
selection, Hopf coefficients and perturbations, using text extraction and
visual checks of critical equations. Exact titles and DOIs, circulant Hopf
reductions, inverse-matrix decay, and five-element enzymatic antecedents
were searched in publisher and author repositories. This is a bounded
literature audit, not an exhaustive priority search. ``Not located'' is
restricted to the inspected source; it does not establish nonexistence in
the literature. Source numerical reports remain reports, even when placed
under a theorem heading.

\subsection{Novelty vocabulary}
The following categories describe the \emph{scope of a claim}, not its proof
status in this project.
\begin{description}[style=nextline]
\item[Definitely known (K).] The stated general fact, model component or
specific antecedent is explicitly present in a checked source.
\item[Directly implied by known theory (D).] Once the project's required
hypotheses are independently verified, a known result supplies the claimed
mechanism or representation. This is an audit inference, not a proof here.
\item[Apparently new specialization (S).] No exact instance was located,
but the proposed contribution applies established machinery to this
particular model or parameter family. Priority remains provisional.
\item[Apparently new theorem (T).] A precisely stated mathematical addition
beyond a routine specialization appears absent from the reviewed literature.
This label requires an identifiable additional conclusion and adequate
literature coverage; it does not mean the theorem is true.
\item[Unresolved because more literature is needed (U).] Source access,
the extent of existing theory or the exact proposed statement is insufficient
for a defensible novelty verdict.
\end{description}
No complete project theorem receives an unconditional T label in this
phase. Quantitative sign regions and parameter-dependent robustness
certificates remain possible theorem contributions, but their exact scope
and nearest antecedents are not sufficiently settled to award that label.
The matrix distinguishes these prospects from established methods.

\section{Six supplied primary sources}
\label{sec:core-sources}

\subsection{S01: Jaramillo, Mrad and Stepien (2023)}
\label{src:jms}
\textbf{Source/version.} \cite{JaramilloMradStepien2023}; supplied
\texttt{2210.04342v2.pdf}, arXiv v2, 14 pages.

\textbf{Exact model and assumptions.} Section III, p.~3, Eqs.~(5)--(6c)
adds global linear mutation to a three-species May--Leonard competition
system. Its first normalized equation is
\[
 \dot m_1=m_1(1-m_1-\alpha m_2-\beta m_3)
            +\mu(-2m_1+m_2+m_3),
\]
with the other two equations cyclically permuted and
$\alpha,\beta,\mu>0$. It has polynomial competition and linear mutation,
not the DM receiver-denominator coupling.

\textbf{Relevant locations.} Section IV.B, p.~5, Theorem 1 and
Eqs.~(7)--(8b): circulant $3\times3$ Jacobian and explicit spectrum.
Section IV.C, pp.~6--7: Theorem 2; general Hopf formula (15), pairing (16),
quadratic tensor (18), critical parameter (19), Fourier vectors
(22)--(23), and explicit coefficient (24)--(25). Equation (20) and the
discussion after (25) exclude the zero-frequency case $\alpha=\beta$.
Section IV.A, pp.~4--5 and Fig.~4, and Section IV.E, pp.~10--11, concern
folds and other coexistence equilibria; these are not the DM bistability
classification.

\textbf{Overlap and verdict.} K: an explicit Fourier-based Hopf calculation
in cyclic competition. D/S: a precedent for P03, P11--P13, not an
arbitrary-$(n,q)$ DM result. None of the eight exact DM targets was located.
The paper itself uses $n$ in a generic normal-form formula on p.~6;
that notation does not generalize its three-species model. Its continuation
and region descriptions have not been independently validated here.

\subsection{S02: Pais, Caicedo-N\'{u}\~{n}ez and Leonard (2012)}
\label{src:pais}
\textbf{Source/version.} \cite{PaisCaicedoLeonard2012}; supplied
\texttt{PaisSIADS12.pdf}, author manuscript pp.~1--30. The published article
is SIADS 11(4), 1754--1784; the pagination is different.

\textbf{Exact model and assumptions.} Section 2, pp.~3--5, Eq.~(2.2):
$\dot x_i=\sum_j x_j f_j(x)q_{ji}-x_i\phi(x)$, with $f=Bx$,
$\phi=x^TBx$, on the simplex $x_i\geq0$, $\sum_i x_i=1$.
Equation (2.1), Condition 1 and mutation rules (Q1), (Q2) specify the
payoff and mutation restrictions. Section 4, p.~10, Eq.~(4.1) uses
$B_{C,N}=\operatorname{Circulant}(1,\alpha,0,\ldots,0,\beta)$,
$\alpha,\beta\in[0,1)$, $\alpha+\beta>0$, $N\geq3$.

\textbf{Relevant locations.} Lemma 4.1, pp.~10--11: mixed equilibrium;
Section 4.1, pp.~11--13, Eqs.~(4.2)--(4.4), Lemmas 4.2 and 4.4,
Theorem 4.5: circulant spectrum, complex pairs, thresholds and crossings.
Section 4.2, p.~13, Lemma 4.6: analytical first Lyapunov coefficient.
Appendix B, pp.~24--25: general tensor/resolvent formula and normalization.
Appendix D, pp.~26--28, Eqs.~(D.1)--(D.4): Fourier reduction, including
the second-harmonic denominator. Section 5, pp.~16--20, treats cycles
with different step sizes; Section 6, pp.~20--23, studies noncirculant
examples numerically.

\textbf{Overlap and verdict.} Strong K methodological overlap with
P04 and P11--P14: directed cyclic networks, arbitrary dimension and
analytical Hopf coefficients are established research territory. The
simplex constraint, mutation mechanism and payoff family differ from DM.
This is not the canonical arbitrary-$(n,q)$ spectrum, DM quadratic
selection theorem or matched DM--MB comparison. Noncirculant simulations
do not supply P15--P18's proposed response or robustness theorems.

\subsection{S03: van Kekem and Sterk (2018)}
\label{src:vks}
\textbf{Source/version.} \cite{vanKekemSterk2018}; supplied
\texttt{1704.05442v3.pdf}, accepted manuscript dated 21 December 2017,
pp.~1--55. Published journal pages are 38--60.

\textbf{Exact model and assumptions.} Section 1.1, p.~1,
Eqs.~(1a)--(1b):
\[
 \dot x_j=x_{j-1}(x_{j+1}-x_{j-2})-x_j+F,
 \qquad x_{j+n}=x_j.
\]
The principal first-instability program assumes $n\geq4$ and increases
$F>0$ from the stable homogeneous state. This quadratic, additively forced
system is distinct from both DM and MB.

\textbf{Relevant locations.} Section 2.1, pp.~8--9, Eqs.~(3)--(6):
Fourier spectrum/vectors; its real symbol uses
$f(j,n)=\cos(2\pi j/n)-\cos(4\pi j/n)$.
Section 2.2, pp.~10--12: Theorems 2 and 4 and Proposition 5, including
simultaneous crossings and the first selected pair. Proposition 5 has
a stated $n=7$ exception and the following remark discusses an $n=12$
tie. Section 2.3, pp.~12--13, Eq.~(8) and Proposition 6: travelling-wave
interpretation and the explicit selected-wave-number limit
$n/l_1(n)\to2\pi/\arccos(1/4)$.
Section 4.3, pp.~38--46, especially (23)--(25) on pp.~39--40 and
(36) on p.~44: general Hopf formula, quadratic reduction and sign analysis.
Section 4.4, pp.~47--48, supplies the proof of Proposition 5.

\textbf{Overlap and verdict.} This is the closest core antecedent for
targets 2--4 and 6. Quadratic-cosine selection, a discrete first-mode
criterion, its large-$n$ limit and a Fourier-reduced Lyapunov coefficient
are already present in another cyclic model. A DM coefficient family,
all its ties/endpoints and its admissible parameter regions were not
located. Their likely category is S rather than a new general Fourier
or selection method. The source's positive-$F$ path must not be silently
identified with a path in DM's five positive parameters.

\subsection{S04: Wei and coauthors (journal issue 2021)}
\label{src:wei}
\textbf{Source/version.} \cite{WeiEtAl2021}; supplied 18-page author
manuscript, filename
\nolinkurl{Wei_et_al_2020_Hopf_bifurcation_and_synchronization.pdf}.
The filename's year and 2020 DOI date do not change the 2021 journal issue.

\textbf{Exact model and assumptions.} Section 2, p.~3, Eq.~(2.1):
\[
\begin{aligned}
 \dot x&=r(y-x)+w,& \dot y&=-(1+m)y+xz-v,\\
 \dot z&=g[1+mx^2-(1+m)xy],&
 \dot w&=2(1+m)w+xz-k_1x,\\
 \dot v&=-mv+k_2y.
\end{aligned}
\]
Here $r,m,g>0$ and $k_1,k_2$ are control parameters. Five is the
dimension of a homopolar dynamo, not the number of identical cyclic
nodes. It is neither a Wu-Xing coupling model nor DM/MB.

\textbf{Relevant locations.} Section 2, pp.~3--5, Eqs.~(2.2)--(2.18),
Theorems 2.1--2.2: a Hopf calculation with rounded parameter/eigenvector
data. Section 3, pp.~5--7, Eqs.~(3.19)--(3.37), Theorems 3.1--3.2:
disturbance observer and controlled convergence. The latter assumes the
specified control and observer construction; it is not passive network
robustness under small heterogeneity. Section 4, pp.~8--9 and the appended
figures supplies control simulations, not a reproduced computation here.

\textbf{Source cautions.} Visually checked p.~3, Eq.~(2.6), describes
five-component vectors but sums only three components in the pairing;
Eq.~(2.7) also has inconsistent displayed dimensions. On p.~4,
Theorem 2.1 both fixes $m=0.04$ and lets $m$ pass a different critical
value. Theorem 2.2 combines a negative reported coefficient with an
unstable-cycle conclusion. These are unresolved manuscript inconsistencies,
not corrected theorems in this audit. Equation (2.10)'s coefficient
normalization also requires convention matching before comparison.

\textbf{Overlap and verdict.} A generic Hopf/tensor calculation is known;
none of the eight exact project targets is located. This paper should
not be cited for a cyclic Fourier reduction, a five-element model or
DM robustness bounds. Check the final publisher text or an erratum
before relying on its detailed criticality claims.

\subsection{S05: Lin and Han (2020)}
\label{src:lin}
\textbf{Source/version.} \cite{LinHan2020}; supplied five-page proceedings
PDF, printed pp.~777--781. Equations (1)--(2), p.~778, were visually
checked because extracted text misreads some subscripts.

\textbf{Exact model and assumptions.} Section 2.1, p.~778, Eq.~(1)
uses six node-dependent coefficient families and additive interventions
$e_i$. Its homogeneous, intervention-free Eq.~(2), written in cyclic
notation for comparison, is
\[
 \dot z_i=a_Lz_i-b_Lz_i^2+c_Lz_iz_{i-1}
       -d_Lz_iz_{i-2}-g_Lz_iz_{i+1}-h_Lz_iz_{i+2},
       \qquad i\in\mathbb Z_5.
\]
The coefficients are explicitly nonnegative. The exact parameter
identification with the canonical five-node MB comparator is
\[
 (a_L,b_L,c_L,d_L,g_L,h_L)=(a,k,c,h,0,g).
\]
This is a comparison of the displayed equations, not a dynamical theorem.
The allowed zero value matters: it makes MB a subfamily, not merely a
limiting resemblance.

\textbf{Relevant locations.} Section 2.2, p.~778: positivity discussion;
Section 2.2.1, pp.~778--779, Cases 1--4: equilibrium examples, numerical
eigenvalues and synchronous scalar reduction. Case 4, p.~779,
Eqs.~(5)--(7), explicitly treats equal initial activities. Sections
2.2.2--2.2.3, pp.~779--781, and Figs.~3--5 report periodic/chaotic
examples and computed Lyapunov exponents. These computations were
not rerun or endorsed here. Section 1 and references [9], [10], [14],
p.~781, point to earlier Wu-Xing models and perturbation studies.

\textbf{Overlap and verdict.} K: the five-node topology, the six-term
polynomial model, its synchronous reduction and MB$_{5,2}$ as a
coefficient subfamily. S/U: the proposed exact comparison between DM
and that comparator, not present here. No DM rational law, arbitrary-$n$
spectrum, q=2 all-$n$ selection theorem or Fourier Green function was
located. Publication of the five-node topology does not establish the
novelty or correctness of the proposed DM analysis.

\subsection{S06: Nian and Li (2025)}
\label{src:nian}
\textbf{Source/version.} \cite{NianLi2025}; supplied publisher PDF
\texttt{s13662-025-03909-9.pdf}, 15 pages, article 61.

\textbf{Exact model and assumptions.} Section 2.2, pp.~5--6,
Eq.~(1), in the source's order $(M,H,T,S,J)$, is
\[
\begin{aligned}
 \dot M&=a_1M+b_1S+c_1J+d_1HTSJ,\\
 \dot H&=a_2H+b_2M+c_2S+d_2MTSJ,\\
 \dot T&=a_3T+b_3H+c_3M+d_3MHSJ,\\
 \dot S&=a_4S+b_4J+c_4T+d_4MHTJ,\\
 \dot J&=a_5J+b_5T+c_5H+d_5MHTS.
\end{aligned}
\]
It has node-dependent linear and quartic interactions. Table 1, p.~7,
gives parameter meanings; Table 4, p.~12, contains fitted coefficients
of mixed signs. The canonical homogeneous positive-parameter domain
must not be imposed on this source.

\textbf{Relevant locations.} Section 2.3, pp.~7--8, Eqs.~(2)--(5):
equilibrium and Jacobian discussion, not a circulant all-$n$ result.
Section 3, pp.~8--10, Eqs.~(6)--(9): parameter fitting.
Section 4, pp.~10--13, Figs.~6--11 and Table 4: illustrative dynamical
and Lyapunov-exponent reports. No exact Hopf coefficient or
heterogeneous Fourier response theorem is given in these sections.

\textbf{Overlap and verdict.} K: an existing nonlinear five-element
application. None of the eight exact mathematical targets was located.
It is not DM, not MB and not evidence that five-element dynamics or
their numerical bifurcation study is new. Its fitted data, chaos
interpretations and displayed Jacobian are not independently validated
by this audit and provide no evidence for our model's biological validity.

\section{Targeted methodological and five-element supplements}
\label{sec:supplemental-sources}
These entries give the closest located general theory or a precise
access gap. Bibliographic inclusion does not imply full-text access.

\subsection{S07: Gray (2006): circulants and their inverses}
\label{src:gray}
\cite{Gray2006}, Stanford author PDF. Section 3.1, printed pp.~32--34
(PDF positions 40--42), Eqs.~(3.7)--(3.13), gives unitary Fourier
diagonalization and explicitly notes normality. Section 3.2,
Theorem 3.1, printed pp.~34--35 (PDF positions 42--43), includes
products and the inverse when all eigenvalues are nonzero.
No specific nonlinear network is assumed. These are K general results;
the DM spectrum and the finite Fourier inverse kernel are D once its
Jacobian and invertibility are verified. Matrix exponentials and
resolvents in the same basis are likewise applications of this
diagonalization, not a new Fourier method. Gray uses a different
Fourier sign convention; indices must be translated. The inverse
formula alone supplies no spatial localization estimate.

\subsection{S08: Kerin and Engler (2020): general quadratic reduction}
\label{src:kerin}
\cite{KerinEngler2020}, arXiv v2, 31-page PDF. Section 3.2,
pp.~11--12, Eqs.~(27)--(29), Lemma 1: an explicit first Lyapunov
coefficient for a real circulant linear part and a quadratic
equivariant nonlinearity. Its formula resolves the uniform and
second-harmonic interactions through the Laurent symbols.
Appendix A, pp.~26--27, Proposition 4 and Eq.~(48), supplies the
bilinear Fourier interaction rule. Section 5, pp.~21--24,
Eq.~(40) and Propositions 1--3 treats site-dependent forcing,
dissipation and advection, assuming the stated energy-preserving
quadratic structure (and positivity of specified diagonal matrices).
K: the quadratic Fourier reduction and heterogeneous-network
precedents. DM is rational and generally has a nonzero third
derivative tensor; its coefficient and sign are not furnished by
Lemma 1. Nor are the source's energy hypotheses established for DM.

\subsection{S09: Stewart (2024): symmetry and mode ordering}
\label{src:stewart}
\cite{Stewart2024}, arXiv v1 HTML. Section 3, Theorem 3.7,
states equivariant Hopf; Section 4.1, Eqs.~(12)--(14), treats a
nearest-neighbor unidirectional ring; Section 5.2, Eq.~(16),
treats longer-range circulant linearizations. Section 5.3,
Theorem 5.5 and Corollary 5.6, permits eigenvalue ordering with
freely chosen connection strengths. Section 5.4, Theorem 5.7,
discusses rotation direction. K: symmetry-based rotating waves
and unrestricted circulant mode control. The sign-constrained
sparse DM family is not the unrestricted interpolation family.
Use its representation-theoretic framework, not genericity as
a replacement for P11--P13. The v1 phase displays in Section 5
have apparent indexing inconsistencies; reconcile them against
the group action before using a time-shift formula.

\subsection{S10: Embree and Trefethen (2001): perturbation bounds}
\label{src:embree}
\cite{EmbreeTrefethen2001}, publisher-form author PDF, pp.~583--590.
Section 1, pp.~583--584, defines pseudospectra via norm-bounded
perturbations and resolvents. Theorem $4_\varepsilon$, p.~585,
gives the similarity/Bauer--Fike inclusion; Theorem 6 and its
discussion, p.~586, address normal matrices. K: general matrix
perturbation machinery. D: a spectral-margin sufficient bound
for an additive matrix perturbation of a stable normal matrix.
An explicit bound on changes of DM parameters must additionally
control the moving equilibrium and the resulting full Jacobian.
This source does not give that parameter certificate or make
the heterogeneous Jacobian normal.

\subsection{S11: Greenbaum, Li and Overton (2020): sensitivity}
\label{src:glo}
\cite{GreenbaumLiOverton2020}, 30-page arXiv author PDF. Section 2,
p.~4, Assumption 1, Theorem 1 and Eq.~(1), gives the derivative
of a simple eigenvalue of an analytic matrix family using
normalized left and right eigenvectors. Section 3.2, p.~12,
states the analytic implicit function theorem used in its
argument. K: simple-eigenvalue sensitivity and the implicit
function framework. D: first-order threshold sensitivity after
the project's parameter family, equilibrium derivative and
nonzero crossing derivative are verified. Multiple eigenvalues
and mode ties require separate treatment. No DM parameter
derivative or explicit robustness radius is supplied.

\subsection{S12: Demko, Moss and Smith (1984): localization caution}
\label{src:demko}
\cite{DemkoMossSmith1984}; partial access only. The publisher-indexed
abstract on p.~491 and Section 2, p.~493, Proposition 2.3,
Eq.~(2.4), and Theorem 2.4 were inspected. They concern decay
of inverses of bounded, boundedly invertible band matrices,
with spectral/conditioning assumptions and stronger forms for
positive definite matrices. K: inverse localization is an
established subject. U: the exact reduction to cyclic graph
distance, wraparound, arbitrary $q$ and uniform-in-$n$ DM
constants. The whole paper and subsequent sparse-matrix
extensions require review before any new localization theorem
is claimed. A finite Fourier sum is not itself a localization
theorem.

\subsection{S13: Lin, Shi and Han (2021): unresolved coupling antecedent}
\label{src:lin2021}
\cite{LinShiHan2021}. The CCC 2021 proceedings contents identify
the title and authors at starting p.~6410; DOI metadata identifies
the article. Full text was not available in the audited corpus.
Thus its exact coupling equations, assumptions and theorem
locations are \emph{unresolved}, not absent. Its enzymatic
five-element topic is directly relevant to targets 1 and 5
and the $n=5,q=2$ motivation. Retrieve and compare the
receiver dependence and any rescaling before claiming the
demand-modulated coupling law is unprecedented.

\subsection{S14: Lin and Han (2024): integrated five-node overlap risk}
\label{src:lin2024}
\cite{LinHan2024}, publisher landing page, abstract
(opening article page 7477) and reference list. The abstract
describes enzymatic regulatory networks with $n\leq9$,
oscillation/adaptation comparisons and a five-node topology
identified with the yin--yang five-element network.
The full equations, analyses and supplement were not accessible.
K: the published topic and claimed topology/application
overlap; U: exact DM law, analytical spectrum and all eight
specific theorem targets. The abstract cannot certify either
equation identity or absence. This is a high-priority main-paper
literature gap, not postponed archival work.

\section{Audit consequences and unresolved evidence}
\label{sec:audit-consequences}

\subsection{What the sources do and do not establish}
No inspected core paper displays $(\DM)_{n,q}$, its
receiver-denominator creation term, or the complete DM--MB
comparison. That is a finite-corpus finding, not a priority
certificate. The five-node bilinear comparator is explicitly
contained in S05. General circulant diagonalization/normality,
Fourier Hopf reduction, rotating waves, inverse response and
matrix/eigenvalue sensitivity all have close antecedents.
The detailed classifications for the eight requested targets
are in Section~\ref{sec:novelty-matrix}.

The paper's possible mathematical contribution should therefore be
judged by the \emph{verified combination of model-specific conclusions}:
the admissible parameter regions, exact mode ties and exclusions,
a calibrated bilinear comparison, a nondegenerate Hopf coefficient
and parameter-dependent response bounds. This audit does not establish
any of them. A change of coupling law does not by itself make every
standard calculation a new theorem.

\subsection{Correction of the initialization's adjoint concern}
S01, p.~6, Eq.~(16), S02, Appendix B, pp.~24--25, and S08,
p.~12, Eq.~(28), use the Hermitian pairing that is conjugate-linear
in its first argument, with
\[
 Jv=\ii\omega v,\qquad J^*p=-\ii\omega p,
 \qquad \langle p,v\rangle=1,
\]
for a real Jacobian ($J^*=J^T$). The opposite adjoint sign is a
standard consistent convention. The initialization's suggestion
that this sign was intrinsically inconsistent is withdrawn in
B07 and claim C078. Phase, eigenvector scale, time scaling,
Taylor factors and the final DM coefficient still require P13
verification. No master text or proof fragment is changed.

\subsection{Priority and access gaps}
\begin{enumerate}
\item Obtain S13 and S14 with their supplementary equations. Compare
the exact nonlinear law, allowable coefficients, variable changes
and normalization; abstracts cannot resolve model priority.
\item Follow S05's p.~781 references [9], [10], [14]: T.~Wang's
2008 Fudan thesis on the dynamic Wu-Xing model; G.~Ding and
T.~Wang's 2008 mathematical analysis, pp.~269--270; and F.~Li's
1987 Yin-yang/Wuxing equations, pp.~966--969. These are
\emph{leads quoted from S05's bibliography}, not independently
audited sources. Their equations and exact result locations are
not asserted here. Resolve them within the main novelty audit.
\item Read the complete S12 paper and follow subsequent inverse-decay
and graph-distance results. Determine whether the planned
localization estimate is a known specialization.
\item Check the broader cyclic-feedback and structured stability-radius
literature once the precise perturbation family, norm and desired
bound are fixed. Current matrix theory settles the general
mechanism, not the novelty of an unspecified final inequality.
\item Preserve edition-specific citations. The final published Wei
text and any corrections have not been compared with the supplied
manuscript; source inconsistencies remain disclosed. The Jaramillo
article number remains a bibliographic completion item.
\end{enumerate}

\subsection{Evidence gates for later work}
P04 must verify the actual DM Jacobian before using circulant
theory, including coincident offsets. P05 must distinguish a
snapshot maximizer from the first crossing on a parameter path;
P06 must specify the sequence of coefficients/branches.
P08 must fix comparator calibration before any superiority
claim. P11--P14 require actual simplicity, nonzero frequency,
transversality, nondegeneracy and the correct group action.
P15--P18 must distinguish a Fourier inverse at the homogeneous
base from the generally noncirculant perturbed problem.
These are audit requirements only. All proof and numerical
files retain their initial TODO/plan status.
